\section{Diamond's two-period economy}\label{sec:diamond}

The base case of the programme, and the one place in it where uniqueness is proved outright rather
than reduced to an assumption. Agents live two periods, work when young and consume out of savings
when old; a representative firm rents capital from the old.

\subsection{The household, in closed form}

A young household with wage $w$ facing gross return $R$ solves $\max u(w-s)+\beta u(Rs)$. With
CRRA utility the Euler equation $u'(w-s)=\beta R\,u'(Rs)$ has the explicit solution $s=\sigma(R)w$
with
\[
  \sigma(R)=\frac{1}{1+\beta^{-1/\gamma}R^{(\gamma-1)/\gamma}}
\]
(\lean{saveShare}). That $\sigma$ satisfies the Euler equation is algebra
(\lean{saveShare\_euler}); that it is the optimum needs one more ingredient.

\begin{proposition}[CRRA lies below its tangent]\label{prop:tangent}
For $c,c'>0$ and $\gamma>0$, $u(c')\le u(c)+u'(c)(c'-c)$.
\formal{crra\_le\_tangent}
\end{proposition}

This is Bernoulli's inequality in both directions. Below unit risk aversion the exponent $1-\gamma$
lies in $(0,1)$ and the concave form applies directly; above it the exponent is negative, and
inverting the consumption ratio and multiplying through turns the claim into the form with
exponent $\gamma\ge1$. Optimality then follows by adding the tangent inequality at the two
consumption levels and cancelling with the Euler equation, with no differentiation of a value
function anywhere (\lean{saveShare\_isOptimal}).

\subsection{The comparative static is discouraging}

\begin{proposition}[The textbook comparative static]\label{prop:diamondcs}
The saving rate rises with the gross return if $\gamma<1$ and falls if $\gamma>1$.
\formal{saveShare\_strictMonoOn}, \lean{saveShare\_strictAntiOn}
\end{proposition}

So the household-side monotonicity that Section~\ref{sec:thm1} cannot establish in general is
here a one-line calculation, and it runs the \emph{wrong} way at the risk aversions economists
use: above unit intertemporal elasticity the income effect beats the substitution effect and a
household facing a better return saves a smaller share of a given wage.

\subsection{The steady state is unique anyway}

What rescues it is general equilibrium. With Cobb-Douglas production the wage rises and the return
falls as capital deepens, and the wage channel dominates. Writing $k'=\sigma(R(k))\,w(k)$ for
capital per worker next period (\lean{diamondMap}), the whole argument is that $k'/k$ falls.

\begin{proposition}[The engine]\label{prop:engine}
If $\theta\le1$, $B>0$, $a\ge0$ and $b>0$, then $x\mapsto x/(1+B(a+bx)^\theta)$ is strictly
increasing on $(0,\infty)$.
\formal{div\_one\_add\_rpow\_strictMono}
\end{proposition}

The proof is the chain $x_2/x_1\le R_2/R_1\le(R_2/R_1)^\theta$, where the first step is
$ax_2\le ax_1$ and the second is that a ratio at most one rises when raised to a power at most one.
No calculus. Substituting $x=k^{\alpha-1}$, which falls with $k$, and $\theta=(\gamma-1)/\gamma$,
which is below one for every positive $\gamma$, gives the result.

\begin{theorem}[Existence and uniqueness]\label{thm:diamond}
For every $\gamma>0$, every $\delta\le1$ and Cobb-Douglas production, the map $k\mapsto k'/k$ is
strictly decreasing, so there is at most one positive steady state; and if the map crosses the
diagonal between two positive levels there is exactly one.
\formal{exists\_unique\_steadyState}, \lean{steadyState\_unique}, \lean{exists\_steadyState}
\end{theorem}

No condition on the savings function is needed, of the kind textbook treatments impose. The
numerical check is in \texttt{numerics/dcheck.m}: the closed form matches a grid search to within
the grid spacing at every risk aversion and return tried; there is exactly one positive crossing
of the diagonal for risk aversions from $0.5$ to $10$ at both $\delta=1$ and $\delta=1/2$; and the
saving rate rises with the return only below unit risk aversion, as
Proposition~\ref{prop:diamondcs} says.

What this does and does not settle is worth stating. It settles uniqueness in the two-period
economy completely, for every risk aversion, which the stochastic life cycle of
Section~\ref{sec:eqm} does not. It settles it by an argument that has no counterpart there,
because with more than two periods the household's saving is not a closed form and $k'/k$ is not a
product of two monotone factors. What it does supply is the shape of the answer: the
household-side comparative static can fail while the general-equilibrium one holds, so a proof of
uniqueness need not go through Light's Theorem 1 at all.
